
\section{Graphical Linear Programming}\label{sec:glp}

With the semantics in place, we define the Symmetric Monoidal Theory (SMT)
that expresses it. In this section we present the base generators and
the new equational setting of the algebra. Further results, such as normal form,
completeness and and interesting examples, are demonstrated
at the end of the section, providing some insightful applications of our theory.

Linear programs combine both the inequality structure
of polyhedra and the notion of an objective function,
and their value functions are precisely the polyhedral convex functions.
The central result of this section is the presentation of a
complete monoidal theory for right-hand side parametric linear programming.

The algebra $\mathbf{GLP}$ extends
$\mathbf{GPA}$ with a single new generator whose semantics encodes the
linear objective function $c^\top x$. As will become clearer later, this
new generator plays a role similar to that of a ``measurement'': it simply exposes
what is in the wire without modifying it.
It lifts $\mathbf{GPA}$ from polyhedral constraints
to full linear programs by attaching a
linear cost to the feasible region.\footnote{An analogous notion was developed by
\cite{stein2024gqa} with their gray arrow generator, introducing the $\ell_2$ norm.}

\begin{definition}[$\mathbf{GLP}$]
	We call $\mathbf{GLP}$ the symmetric monoidal theory obtained by
	extending $\mathbf{GPA}$ with the following generator:
	\begin{figure}[H]
		\centering
            \begin{tikzpicture}
    \node at (0.2,0.5) {$\llbracket$};
    \pic (g1) at (0.500,0.500) {linear};
    \node at (1.25,0.5) {$\rrbracket$};
    \node at (1.65,0.5) {$=$};
    \node[anchor=west] at (1.9,0.5) {$(\bullet,x) \mapsto x$};
\end{tikzpicture}

	\end{figure}
	together with the following axioms governing its interaction with
	the existing generators:

    \begin{tabular}{r c c c}
      \diagrow{(linearity)}{ 
          \begin{scope}[shift={(0.000,0.500)}]
          \pic (g1) at (0.500,0.500) {linear};
          \end{scope}
          \pic (g2) at (0.500,0.500) {linear};
      }{ 
          \pic (g1) at (0.500,0.500) {linear};
          \begin{scope}[shift={(1.250,0.000)}]
          \pic (g2) at (0.500,0.500) {comultiplication};
          \end{scope}
          \draw (g1-out-0) to[out=0,in=180] (g2-in-0);
      }
      \diagrow{(zeroeq)}{ 
          \pic (g1) at (0.500,0.500) {linear};
          \begin{scope}[shift={(1.250,0.000)}]
          \pic (g2) at (0.500,0.500) {cozero};
          \end{scope}
          \draw (g1-out-0) to[out=0,in=180] (g2-in-0);
      }{ 
          \node[draw, dotted] {};
      }
      \diagrow{(inf)}{
          \pic (g1) at (0.500,0.500) {linear};
          \begin{scope}[shift={(1.250,0.000)}]
          \pic (g2) at (0.500,0.500) {ge};
          \end{scope}
          \draw (g1-out-0) to[out=0,in=180] (g2-in-0);
      }{
          \pic (g1) at (0.500,0.500) {linear};
      }
    \end{tabular}
\end{definition}

The axioms represent expected properties of the linear cost function.
The \textit{linearity} axiom expresses $f(x) + f(y) = f(x+y)$, where $f$ denotes the interpretation of $\inlinegen{linear}$ and the sum denotes vertical composition; composing two copies of $\inlinegen{linear}$ vertically is the same as applying a single $\inlinegen{linear}$ to the combined wires. Furthermore, this axiom is also the fundamental property of LPs, which is the fact that they are invariant to column stochastic matrices (see~\ref{rmk:stochastic}).

The \textit{inf} axiom expresses the defining property of minimization: the infimum of a quantity over all values satisfying a lower bound is that lower bound itself,
\[\inf_x \{ x \mid x \ge y \} =  y,\]
so applying $\inlinegen{ge}$ after $\inlinegen{linear}$ has no effect.
Finally, the \textit{zeroeq} axiom captures the tautological case $\inf\{x \mid x = 0\} = 0$: the minimum of a linear cost subject to the equality $x = 0$ is zero.

\begin{definition}
	We call $\mathrm{FreeP}_{\mathbf{GLP}}$ the PROP freely generated
	by the signature and equations of $\mathbf{GLP}$.
\end{definition}

The normal form theorem for $\mathbf{GLP}$ extends the one for
$\mathbf{GPA}$: every diagram can be reduced to a $\mathbf{GPA}$ block
encoding the constraints, composed with the new generator encoding the
objective. The proof proceeds by structural induction, using the existing
normal form for $\mathbf{GPA}$ as the inductive base; it is given in
Appendix~\ref{app:proofs}.

\begin{theorem}[Normal Form]\label{thm:NormalFormGLP}
	For every diagram $c \in \mathbf{GLP}$, there exists a diagram $c'$
	in normal form such that $c = c'$:
	\begin{figure}[H]
		\centering
        \begin{tikzpicture}
    \pic (g1) at (0.500,1.300) {linear};
    \coordinate (id_in_2) at (0,0.700);
    \coordinate (id_out_3) at (1.000,0.700);
    \draw (id_in_2) -- (id_out_3);
    \pic[val=A] (g4) at (2.000,1.000) {2drelation};
    \draw (g1-out-0) to[out=0,in=180] (g4-in-0);
    \draw (id_out_3) to[out=0,in=180] (g4-in-1);
    \pic (g5) at (3.500,1.000) {ge};
    \draw (g4-out-0) to[out=0,in=180] (g5-in-0);
    \pic[val=B] (g6) at (5.000,1.000) {corelation};
    \draw (g5-out-0) to[out=0,in=180] (g6-in-0);
\end{tikzpicture}

	\end{figure}
\end{theorem}

The following corollary reveals the geometric meaning of the normal form:
the $\mathbf{GPA}$ part of every $\mathbf{GLP}$ diagram encodes the
epigraph of the corresponding value function. This is the key structural
insight connecting the syntax to the semantic equivalence of
Proposition~\ref{prop:equalepi}.

\begin{corollary}[Epigraph Form]\label{cor:epigraph form}
	Given a diagram $d \in \mathbf{GLP}$ in normal form, it can always
	be rewritten as a diagram $d'$ of the form:
	\begin{figure}[H]
		\centering
		\begin{tikzpicture}
    \pic (g1) at (0.500,0.500) {linear};
    \begin{scope}[shift={(1.250,0.000)}]
        \pic (g2) at (0.500,0.500) {ge};
    \end{scope}
    \begin{scope}[shift={(2.500,0.000)}]
        \pic (g3) at (0.500,0.500) {comultiplication};
    \end{scope}
    \begin{scope}[shift={(3.750,0.000)}]
        \pic[val=A'] (g4) at (0.500,0.500) {2dmatrix};
    \end{scope}
    \begin{scope}[shift={(5.000,0.000)}]
        \pic (g5) at (0.500,0.500) {ge};
    \end{scope}
    \begin{scope}[shift={(6.250,0.000)}]
        \pic[val=B] (g6) at (0.500,0.500) {corelation};
    \end{scope}
    \draw (g1-out-0) to[out=0,in=180] (g2-in-0);
    \draw (g2-out-0) to[out=0,in=180] (g3-in-0);
    \draw (g3-out-0) to[out=0,in=180] (g4-in-0);
    \draw (g3-out-1) to[out=0,in=180] (g4-in-1);
    \draw (g4-out-0) to[out=0,in=180] (g5-in-0);
    \draw (g5-out-0) to[out=0,in=180] (g6-in-0);
    \draw[dotted] (1.15,0.10) rectangle (7.05,0.90);
\end{tikzpicture}
	\end{figure}
	where the semantics of the dotted diagram is
	\[\left\{(t,y)\mid \exists\, x,\; A'x \ge By + c,\; t\ge x\; \right\}.\]
	This is the epigraph of the original diagram $d$
	\begin{align*}
		[d](y)            & = \inf_{x} \; x \quad \text{s.t. } A' x \ge B y + c \\
		\mathrm{epi}([d]) & = \left\{(t, y) \mid \exists\, x \le t,\;
		A' x \ge B y + c \right\}
	\end{align*}
	We call this the \emph{epigraph form} of a linear program.
\end{corollary}

\subsection{Soundness and Completeness}\label{sec:complete}

We now show that $\mathbf{GLP}$ is sound, expressive, and complete with
respect to $\mathbf{LPRel}$, establishing that the semantics of parametric
linear programs is fully axiomatised. All proofs are given in
Appendix~\ref{app:proofs}.

\begin{lemma}[Soundness]\label{lem:soundness}
	If $s = t$ in $\mathrm{FreeP}_{\mathbf{GLP}}$, then $[s] = [t]$
	in $\mathbf{LPRel}$.
\end{lemma}

\begin{lemma}[Expressivity]\label{lem:expressivity}
	For every $f : m \to n$ in $\mathbf{LPRel}$ there exists a
	$\Sigma$-term $s : m \to n$ in
	$\mathrm{Hom}_{\mathrm{FreeP}_{\mathbf{GLP}}}(m,n)$
	such that $[s] = f$.
\end{lemma}

The final lemma is completeness. Its proof brings together the main
tools developed throughout: the epigraph form
(Corollary~\ref{cor:epigraph form}), the shared epigraph property
(Proposition~\ref{prop:equalepi}), the completeness of $\mathbf{GPA}$,
and the FME-based normal form.

\begin{lemma}[Completeness]\label{lem:completeness}
	If $[c] = [d]$ in $\mathbf{LPRel}$, then $c = d$ in
	$\mathrm{FreeP}_{\mathbf{GLP}}$.
\end{lemma}

As a consequence of the three lemmas above:

\begin{theorem}\label{thm:main}
	$\mathbf{LPRel}$ is presented by $\mathrm{FreeP}_{\mathbf{GLP}}$.
\end{theorem}
\begin{proof}
	By the preceding three lemmas.
\end{proof}

\subsection{Examples}\label{sec:examples}

The expressivity of $\mathbf{GLP}$ becomes concrete when one considers specific classes of piecewise affine convex functions that arise naturally across mathematics and engineering. Each example below is a morphism in $\mathbf{LPRel}$ and can therefore be represented as a $\mathbf{GLP}$ diagram in normal form.

\begin{example}[ReLU activation]
  The rectified linear unit $f(y) = \max(0, y)$ is the elementary building block of modern neural networks. It is a piecewise affine convex function, and therefore a morphism in $\mathbf{GLP}$, via the parametric linear program
  \[
    f(y) = \min_{x} \{ x \mid x \geq 0,\; x \geq y \}.
  \]
  As a $\mathbf{GLP}$ diagram, it consists of a single $\inlinegen{linear}$ generator preceded by two inequality constraints: one bounding $x$ from below by $y$ and one enforcing non-negativity.

  \begin{center}
    \vspace{1.5em}
    \begin{tikzpicture}
  \begin{scope}[shift={(0.000,0.500)}]
  \pic (g1) at (0.500,0.500) {linear};
  \begin{scope}[shift={(1.250,0.000)}]
  \pic (g2) at (0.500,0.500) {copy};
  \end{scope}
  \draw (g1-out-0) to[out=0,in=180] (g2-in-0);
  \end{scope}
  \begin{scope}[shift={(2.500,0.000)}]
  \begin{scope}[shift={(0.625,1.000)}]
  \pic (g3) at (0.500,0.500) {ge};
  \end{scope}
  \pic (g4) at (0.500,0.500) {ge};
  \begin{scope}[shift={(1.250,0.000)}]
  \pic (g5) at (0.500,0.500) {cozero};
  \end{scope}
  \draw (g4-out-0) to[out=0,in=180] (g5-in-0);
  \coordinate (a6) at (0.000,0 |- g3-in-0);
  \draw (g3-in-0) -- (a6);
  \coordinate (a7) at (2.250,0 |- g3-out-0);
  \draw (g3-out-0) -- (a7);
  \end{scope}
  \draw (g2-out-0) to[out=0,in=180] (a6);
  \draw (g2-out-1) to[out=0,in=180] (g4-in-0);
\end{tikzpicture}

    \vspace{1em}
  \end{center}
\end{example}

\begin{example}[Feedforward networks with ReLU activations]
  A one-hidden-layer feedforward network with non-negative output weights produces a piecewise affine convex function:
  \[
    f(y) = \sum_{i=1}^{k} c_i \max(0, w_i^\top y + b_i), \qquad c_i \geq 0.
  \]
  Each summand $c_i \max(0, w_i^\top y + b_i)$ is piecewise affine convex (a scaled ReLU composed with an affine map), and their sum is again piecewise affine convex. The combined function is the value function of the parametric linear program
  \[
    f(y) = \min_{x_1,\ldots,x_k} \sum_{i=1}^{k} c_i x_i \quad \text{s.t.} \quad x_i \geq 0,\; x_i \geq w_i^\top y + b_i, \quad i = 1,\ldots,k.
  \]
  The compositional structure of the network maps directly onto the monoidal structure of $\mathbf{GLP}$: neurons within the same layer compose in parallel via the tensor product, while successive layers compose sequentially via categorical composition.

  \begin{center}
    \vspace{1.5em}
    \diagdef{
    \node[draw, fill=white] at (0,0) {ReLU};
    \draw (0.55,0.0) to (1.0,0);
    \draw (-0.55,0.0) to (-1.0,0);
  }{
  \begin{scope}[shift={(0.000,0.500)}]
  \begin{scope}[shift={(1.250,0.000)}]
  \pic (g2) at (0.500,0.500) {copy};
  \end{scope}
  \end{scope}
  \begin{scope}[shift={(2.500,0.000)}]
  \begin{scope}[shift={(0.625,1.000)}]
  \pic (g3) at (0.500,0.500) {ge};
  \end{scope}
  \pic (g4) at (0.500,0.500) {ge};
  \begin{scope}[shift={(1.250,0.000)}]
  \pic (g5) at (0.500,0.500) {cozero};
  \end{scope}
  \draw (g4-out-0) to[out=0,in=180] (g5-in-0);
  \coordinate (a6) at (0.000,0 |- g3-in-0);
  \draw (g3-in-0) -- (a6);
  \coordinate (a7) at (2.250,0 |- g3-out-0);
  \draw (g3-out-0) -- (a7);
  \end{scope}
  \draw (g2-out-0) to[out=0,in=180] (a6);
  \draw (g2-out-1) to[out=0,in=180] (g4-in-0);
  }

\end{center}
\begin{center}
    \vspace{1em}
    \begin{tikzpicture}
  \pic (g1) at (0.500,0.500) {linear};
  \begin{scope}[shift={(1.250,0.000)}]
  \pic (g2) at (0.500,0.500) {relu};
  \end{scope}
  \draw (g1-out-0) to[out=0,in=180] (g2-in-0);
  \begin{scope}[shift={(2.500,0.000)}]
      \pic[val=W_1] (g3) at (0.500,0.500) {corelation};
  \end{scope}
  \draw (g2-out-0) to[out=0,in=180] (g3-in-0);
  \begin{scope}[shift={(3.750,0.000)}]
  \pic (g4) at (0.500,0.500) {relu};
  \end{scope}
  \draw (g3-out-0) to[out=0,in=180] (g4-in-0);
  \begin{scope}[shift={(5.000,0.000)}]
      \pic[val=W_2] (g5) at (0.500,0.500) {corelation};
  \end{scope}
  \draw (g4-out-0) to[out=0,in=180] (g5-in-0);
\end{tikzpicture}

  \end{center}
\end{example}

\begin{example}[$\ell_1$ regularisation]
  The $\ell_1$ norm $f(y) = \|y\|_1 = \sum_{i=1}^{n} |y_i|$ is a canonical example in sparse optimisation~\cite{boyd2004convex}. It admits the LP representation
  \[
    f(y) = \min_{t \in \mathbb{R}^n} \mathbf{1}^\top t \quad \text{s.t.} \quad t_i \geq y_i,\; t_i \geq -y_i, \quad i = 1,\ldots,n,
  \]
  making $\|y\|_1$ a $\mathbf{GLP}$ morphism of type $n \to 0$. The absolute value of each component, $|y_i| = \max(y_i, -y_i)$, is itself a piecewise affine convex function, and their sum is obtained via the additive structure inherited from $\mathbf{GLA}$.

  \begin{center}
    \vspace{1.5em}
    \begin{tikzpicture}
  \begin{scope}[shift={(0.000,0.500)}]
  \pic (g1) at (0.500,0.500) {linear};
  \end{scope}
  \begin{scope}[shift={(1.250,0.000)}]
  \begin{scope}[shift={(0.000,0.500)}]
  \pic (g2) at (0.500,0.500) {copy};
  \end{scope}
  \begin{scope}[shift={(1.250,0.000)}]
  \begin{scope}[shift={(0.625,1.000)}]
  \pic (g3) at (0.500,0.500) {ge};
  \end{scope}
  \pic (g4) at (0.500,0.500) {ge};
  \begin{scope}[shift={(1.250,0.000)}]
  \pic (g5) at (0.500,0.500) {antipode};
  \end{scope}
  \draw (g4-out-0) to[out=0,in=180] (g5-in-0);
  \coordinate (a6) at (0.000,0 |- g3-in-0);
  \draw (g3-in-0) -- (a6);
  \coordinate (a7) at (2.250,0 |- g3-out-0);
  \draw (g3-out-0) -- (a7);
  \end{scope}
  \draw (g2-out-0) to[out=0,in=180] (a6);
  \draw (g2-out-1) to[out=0,in=180] (g4-in-0);
  \begin{scope}[shift={(3.750,0.500)}]
  \pic (g8) at (0.500,0.500) {cocopy};
  \end{scope}
  \draw (a7) to[out=0,in=180] (g8-in-0);
  \draw (g5-out-0) to[out=0,in=180] (g8-in-1);
  \end{scope}
  \draw (g1-out-0) to[out=0,in=180] (g2-in-0);
\end{tikzpicture}

    \vspace{1em}
  \end{center}
\end{example}
